\begin{corollary}[局部信息密度的一阶离散修正律]\label{cor:xi-fold-local-information-density-first-order-discrete-law}
令 \(W_m\sim\mu_m\)，并定义
\[
\gamma_m:=\frac{1}{m}\log d_m^{\mathrm{bin}}(W_m).
\]
再令 \(B\) 为参数
\[
\mathbb P(B=1)=\frac{1}{\varphi\sqrt5},
\qquad
\mathbb P(B=0)=\frac{\varphi}{\sqrt5}
\]
的 Bernoulli 变量。则有分布收敛
\[
m\left(\gamma_m-\log\!\frac{2}{\varphi}\right)
\Longrightarrow
-(\log\varphi)B.
\]
并且
\[
\lim_{m\to\infty}
m\left(\mathbb E[\gamma_m]-\log\!\frac{2}{\varphi}\right)
=
-\frac{\log\varphi}{\varphi\sqrt5},
\]
\[
\lim_{m\to\infty}
m^2\operatorname{Var}(\gamma_m)
=
(\log\varphi)^2
\frac{1}{\varphi\sqrt5}
\left(
1-\frac{1}{\varphi\sqrt5}
\right).
\]
因此 \(\gamma_m\) 的首个非平凡修正并不呈现 \(\sqrt m\)-Gaussian 涨落，而是在 \(m^{-1}\) 尺度上塌缩为末位几何所决定的两点律。
\end{corollary}
\begin{proof}
在定理 \ref{thm:xi-fold-escort-log-multiplicity-two-atom} 中取 \(q=1\)。
由于 \(\mu_{m,1}=\mu_m\)，且
\[
X_m^{(1)}
=
\log d_m^{\mathrm{bin}}(W_m)-m\log\!\Bigl(\frac{2}{\varphi}\Bigr)
=
m\left(\gamma_m-\log\!\frac{2}{\varphi}\right),
\]
故该定理立即给出
\[
m\left(\gamma_m-\log\!\frac{2}{\varphi}\right)
\Longrightarrow
\frac{\varphi}{\sqrt5}\delta_0+\frac{1}{\varphi\sqrt5}\delta_{-\log\varphi},
\]
这正是随机变量 \(-(\log\varphi)B\) 的分布。

又由定理 \ref{thm:xi-fold-escort-log-multiplicity-two-atom} 的证明，
\[
X_m^{(1)}=-W_m(m)\log\varphi+O\!\Bigl(\Bigl(\frac{\varphi}{2}\Bigr)^m\Bigr),
\]
且误差一致于样本点，从而 \((X_m^{(1)})_{m\ge 1}\) 一致有界。
因此其一、二阶矩都收敛到极限分布的相应矩。于是
\[
\lim_{m\to\infty}\mathbb E[X_m^{(1)}]
=
\mathbb E[-(\log\varphi)B]
=
-\frac{\log\varphi}{\varphi\sqrt5},
\]
\[
\lim_{m\to\infty}\operatorname{Var}(X_m^{(1)})
=
\operatorname{Var}((\log\varphi)B)
=
(\log\varphi)^2
\frac{1}{\varphi\sqrt5}
\left(
1-\frac{1}{\varphi\sqrt5}
\right).
\]
最后用
\[
X_m^{(1)}=m\left(\gamma_m-\log\!\frac{2}{\varphi}\right)
\]
即可得到所示均值与方差极限。证毕。
\end{proof}
